\documentclass{report}
\usepackage{amsmath,amssymb}
\setlength{\parindent}{0mm}
\setlength{\parskip}{1em}
\begin{document}
\begin{center}
\rule{6in}{1pt} \
{\large Kathryn Lund-Nguyen \\
{\bf The Radau-Lanczos method for matrix functions}}

Department of Mathematics \\ Temple University \\ 1805 N Broad Street \\ Rm 638 Wachman Hall \\ Philadelphia \\ PA 19122
\\
{\tt lund-nguyen@temple.edu}\\
Andreas Frommer\\
Marcel Schweitzer\\
Daniel B. Szyld\end{center}

We present a new iterative method for computing $f(A) \vb$, derived from
a relationship between the standard Lanczos method and a Gauss-Radau
quadrature rule. We show that this method, called the Radau-Lanczos
method, converges when $A$ is Hermitian positive definite and $f$ is a
Stieltjes function. We also show that the restarted version of this
method converges and present numerical results showing this method
performing better than the standard Lanczos method in terms of attainable
error norm and iteration count.


\end{document}
